% -*- coding: utf-8 -*-
\chapter{英語論文の執筆と推敲}
\label{chap4}

\begin{chapterbox}
\textbf{【この章で学ぶこと】}
\begin{itemize}
\item \textbf{英語論文の基本構造（IMRaD）を理解する} --- 各セクションの役割と書くべき内容を把握し、構造的な論文を書ける
\item \textbf{AIを活用した英訳・執筆・推敲ワークフローを習得する} --- 日本語原稿からAIの支援を受けて高品質な英語論文を仕上げるプロセスを実践できる
\item \textbf{学術的な英語表現を適切に使い分けられる} --- 分野に応じた定型表現を活用し、査読者に好印象を与える論文を執筆できる
\end{itemize}
\end{chapterbox}
\vspace{5mm}

%======================================================================
\section{なぜ英語論文が必要か}
\label{sec:4.1}
\index{えいごろんぶん@英語論文}

\subsection{国際会議・ジャーナルの重要性}
\index{こくさいかいぎ@国際会議}
\index{じゃーなる@ジャーナル}

研究の価値は、その成果が広く共有され、他の研究者に活用されて初めて発揮される。日本語で発表された研究は、国内の研究者には届くが、世界中の研究者にはほとんど届かない。トップカンファレンス（NeurIPS、ICML、ACL、CVPR、AAAI等）やトップジャーナル（Nature、Science、IEEE Transactions等）はすべて英語で発表されており、国際的な研究コミュニティの共通言語は英語である。

大学院生にとって英語論文の執筆は、以下の理由で極めて重要である。

\begin{itemize}
\item \textbf{研究成果の国際的な可視性:} 英語論文は世界中の研究者が読むことができ、被引用数の向上につながる
\item \textbf{キャリア形成:} 国際的な実績は、アカデミアでも産業界でも高く評価される
\item \textbf{研究コミュニティへの参加:} 英語論文を発表することで、国際的な研究ネットワークに参加できる
\item \textbf{フィードバックの質:} 国際会議の査読者は世界トップレベルの研究者であり、質の高いフィードバックが得られる
\end{itemize}

\subsection{研究成果の国際発信と評価}

英語論文の執筆は、単に日本語を英語に翻訳する作業ではない。国際的な読者に向けて、研究の背景、意義、手法、結果を論理的かつ説得力のある形で伝える作業である。

しかし、英語を母語としない研究者にとって、学術英語の壁は非常に高い。従来は英文校正サービスやネイティブの共同研究者に頼ることが一般的であったが、AIの登場により、英語論文執筆のハードルは大きく下がった。AIは万能ではないが、適切に活用すれば、執筆の効率と品質を大幅に向上させられる。

%======================================================================
\section{英語論文の基本構造}
\label{sec:4.2}

\subsection{IMRaD構造}
\index{IMRaD@IMRaD構造}

国際的な学術論文の標準的な構造はIMRaDと呼ばれ、以下の4つのセクションから構成される。

\begin{table}[htbp]
\centering
\caption{IMRaD構造の概要}
\label{tab:imrad}
\footnotesize
\begin{tabular}{llll}
\toprule
セクション & 英語名 & 主な内容 & 時制 \\
\midrule
はじめに & Introduction & 研究の背景、目的、貢献 & 現在形・過去形 \\
手法 & Methods & 提案手法の詳細 & 過去形 \\
結果 & Results & 実験結果の提示 & 過去形 \\
考察 & Discussion & 結果の解釈、限界、将来展望 & 現在形・過去形 \\
\bottomrule
\end{tabular}
\end{table}

これに加えて、以下のセクションが一般的に含まれる。

\begin{itemize}
\item \textbf{Abstract（要旨）:} 論文全体の要約（150--300語程度）
\item \textbf{Related Work（関連研究）:} 先行研究の整理と本研究の位置づけ
\item \textbf{Conclusion（結論）:} 研究の要約と将来展望
\item \textbf{References（参考文献）:} 引用した文献のリスト
\end{itemize}

\subsection{各セクションの役割と書くべき内容}

\begin{itembox}[l]{Introduction（はじめに）の構造}
Introductionは「逆三角形」の構造で書く。まず広い背景から始めて、徐々に焦点を絞り、最後に自分の研究の貢献を明確に述べる。
\begin{enumerate}
\item 研究分野の一般的な背景と重要性（1--2段落）
\item 既存研究の課題や未解決の問題（1--2段落）
\item 本研究のアプローチと貢献（1段落）
\item 論文の構成の概要（1段落、省略可）
\end{enumerate}
\end{itembox}

\textbf{Methods（手法）:}

提案手法を、読者が再現できるレベルで詳細に記述する。

\begin{enumerate}
\item 問題設定の形式的な定義
\item 提案手法の全体像
\item 各構成要素の詳細
\item 実装の詳細（ハイパーパラメータ、使用ライブラリ等）
\end{enumerate}

\textbf{Results（結果）:}

実験結果を客観的に提示する。解釈はDiscussionで行う。

\begin{enumerate}
\item 実験設定（データセット、評価指標、ベースライン）
\item 主要な結果（表やグラフ）
\item 結果の記述（数値の比較、傾向の指摘）
\end{enumerate}

\textbf{Discussion（考察）:}

結果の意味を解釈し、研究の意義と限界を議論する。

\begin{enumerate}
\item 主要な発見の解釈
\item 先行研究との比較
\item 研究の限界
\item 将来の研究方向
\end{enumerate}

\subsection{パラグラフの構成法}
\index{ぱらぐらふ@パラグラフ}

英語の学術論文では、パラグラフ（段落）が意味の基本単位である。各パラグラフは以下の構造に従う。

\begin{enumerate}
\item \textbf{トピックセンテンス:} パラグラフの主張を1文で述べる（段落の冒頭に置く）
\item \textbf{サポート:} トピックセンテンスを支える根拠、例、データを述べる（2--5文）
\item \textbf{結論/遷移:} パラグラフの内容をまとめ、次のパラグラフへつなげる（1文）
\end{enumerate}

\begin{promptbox}
以下の日本語の段落を、トピックセンテンス・サポート・結論の構造に
整理してください。必要に応じて、論理の飛躍を指摘してください。

{[}日本語の段落を貼り付け{]}
\end{promptbox}

\begin{itembox}[l]{Tip: 1パラグラフ1メッセージ}
英語論文では「1パラグラフ1メッセージ」の原則を守ることが重要である。複数のメッセージが混在するパラグラフは、読者の理解を妨げる。日本語の文章は1つの段落に複数の話題を含むことがあるが、英語論文ではそれぞれを独立したパラグラフに分けるのが原則である。
\end{itembox}

%======================================================================
\section{AIを使った英訳ワークフロー}
\label{sec:4.3}
\index{えいやくわーくふろー@英訳ワークフロー}

\subsection*{Step 1: 日本語原稿の準備と構造化}

AIに英訳を依頼する前に、日本語原稿の品質を高めることが重要である。曖昧な日本語からは曖昧な英語しか生まれない。

\textbf{日本語原稿の準備チェックリスト:}

\begin{itemize}
\item 各段落にトピックセンテンスがあるか
\item 論理の飛躍がないか
\item 主語と述語が明確か（日本語は主語省略が多いので注意）
\item 専門用語が一貫して使われているか
\item 図表への参照が正しいか
\end{itemize}

\textbf{構造化のためのプロンプト:}

\begin{promptbox}
以下の日本語原稿を、英語論文に翻訳する前の準備として、
構造的に整理してください。

具体的には:
1. 各段落のトピックセンテンスを明示してください
2. 論理の飛躍や不明確な箇所を指摘してください
3. 主語が省略されている文を特定し、主語を補ってください
4. 段落の分割・統合が必要な箇所を提案してください

{[}日本語原稿を貼り付け]
\end{promptbox}

\subsection*{Step 2: セクション単位でAIに英訳を依頼}

日本語原稿が整理できたら、セクション単位でAIに英訳を依頼する。論文全体を一度に訳すのではなく、セクション単位で行うことで、品質の管理がしやすくなる。

\textbf{効果的な英訳プロンプト:}

\begin{promptbox}
以下の日本語原稿を、学術英語に翻訳してください。

【条件】
- 分野: [例: コンピュータサイエンス / 機械学習]
- 対象読者: 国際会議の査読者（専門家）
- 文体: 学術論文（フォーマル）
- セクション: [例: Introduction / Methods / Results]
- 時制: [例: 過去形を基本とする（Methodsの場合）]

【注意事項】
- 以下の専門用語は指定の英語表現を使用してください:
  [用語1]: [英語表現1]
  [用語2]: [英語表現2]
- 受動態を過度に使用しないでください
- 1文が長くなりすぎないようにしてください（目安: 25語以内）

【日本語原稿】
{[}原稿を貼り付け]
\end{promptbox}

\begin{itembox}[l]{分野指定の重要性}
学術英語の表現は分野によって大きく異なる。例えば、同じ「提案する」でも分野によって適切な表現が異なる。
\begin{itemize}
\item コンピュータサイエンス: ``We propose...'' / ``We present...''
\item 医学: ``We hypothesize...'' / ``We examined...''
\item 物理学: ``We derive...'' / ``We demonstrate...''
\end{itemize}
\end{itembox}

\begin{itembox}[l]{Tip: 参考論文の活用}
翻訳の際は、自分の分野のトップ論文を2--3本手元に置いておき、表現や文体の参考にするとよい。AIが生成した英語が、その分野の慣習に合っているかを確認するための基準になる。
\end{itembox}

\subsection*{Step 3: 学術的表現への洗練}

初回の翻訳結果をさらに学術的に洗練する。

\begin{promptbox}
以下の英語論文の草稿を、より学術的で洗練された表現に改善してください。

【改善の方向性】
1. 曖昧な表現をより具体的・精確にする
2. 冗長な表現を簡潔にする
3. 接続語を適切に使い、論理の流れを明確にする
4. 学術論文で一般的な表現に置き換える
   （例: "very good" → "substantially improved"）
5. ヘッジング表現を適切に使用する
   （例: "This proves..." → "These results suggest..."）

改善箇所には【変更理由】を付けてください。

{[}英語草稿を貼り付け]
\end{promptbox}

\subsection*{Step 4: 文法・スタイルチェック}

最終段階として、文法とスタイルの統一性を確認する。

\begin{promptbox}
以下の英語論文について、文法・スタイルの最終チェックを行ってください。

【チェック項目】
1. 文法エラー（主語-動詞の一致、冠詞、前置詞など）
2. 時制の一貫性（セクション内で統一されているか）
3. 用語の一貫性（同じ概念に同じ用語が使われているか）
4. 数値の表記（有効数字、単位の表記）
5. 参照の形式（Figure 1、Table 2などの表記統一）
6. スペルの統一（British/American Englishの混在がないか）

エラーが見つかった場合、修正案と理由を示してください。

{[}英語論文を貼り付け]
\end{promptbox}

%======================================================================
\section{推敲テクニック}
\label{sec:4.4}
\index{すいこう@推敲}

\subsection{段階的推敲法}

英語論文の推敲は、一度に全てをチェックしようとすると、多くの問題を見落とす。以下の3段階に分けて、段階的に推敲を行うとよい。

\textbf{第1段階: 文法チェック}

\begin{promptbox}
以下の英語論文について、文法エラーのみに焦点を当ててチェックしてください。
表現の改善提案は不要です。純粋な文法エラー（主語-動詞の一致、冠詞、
前置詞、時制の誤りなど）のみを指摘し、修正案を示してください。

{[}英語論文を貼り付け]
\end{promptbox}

\textbf{第2段階: 表現チェック}

\begin{promptbox}
以下の英語論文について、表現の改善に焦点を当ててチェックしてください。
文法は正しいが、以下の観点で改善できる箇所を指摘してください。

1. より学術的な表現への置き換え
2. 冗長な表現の簡潔化
3. 曖昧な表現の明確化
4. 不自然な表現（日本語直訳調）の修正

{[}英語論文を貼り付け]
\end{promptbox}

\textbf{第3段階: 論理チェック}

\begin{promptbox}
以下の英語論文について、論理構造に焦点を当ててチェックしてください。

1. 段落間の論理的つながりは適切か
2. 主張に対する根拠は十分か
3. 飛躍や矛盾はないか
4. 読者が混乱する可能性のある箇所はどこか

{[}英語論文を貼り付け]
\end{promptbox}

\subsection{比較推敲法}

複数のAIツールで翻訳を行い、最良の表現を選ぶ方法である。

\begin{promptbox}
以下の日本語を3つの異なるスタイルで英訳してください。

スタイルA: 簡潔で直接的な表現
スタイルB: 詳細で丁寧な表現
スタイルC: [分野名]の国際会議で一般的な表現

それぞれの訳について、長所と短所を説明してください。

{[}日本語原稿を貼り付け]
\end{promptbox}

さらに、Claude、ChatGPT、Geminiなど異なるAIで同じ文章を翻訳し、表現を比較することも効果的である。各AIには得意な表現スタイルがあり、それぞれの良い部分を組み合わせることで、より高品質な英語論文を作成できる。

\subsection{逆翻訳チェック}
\index{ぎゃくほんやくちぇっく@逆翻訳チェック}

英訳が原文の意味を正確に反映しているかを確認する強力な方法が、逆翻訳チェックである。

\begin{promptbox}
以下の英語を日本語に翻訳してください。
できるだけ直訳に近い形で翻訳し、英語の構造がわかるようにしてください。

{[}英訳された論文の一部を貼り付け]
\end{promptbox}

逆翻訳された日本語を元の日本語原稿と比較し、意味のズレがないかを確認する。ズレが見つかった場合、英訳を修正する。

\textbf{逆翻訳で発見できる問題:}

\begin{itemize}
\item 原文にあった情報が英訳で欠落している
\item 英訳が原文とは異なる意味に解釈される
\item 曖昧な英語表現が原文の意図と異なる解釈を生んでいる
\item 専門用語の翻訳が不正確である
\end{itemize}

\subsection{ネイティブチェック風レビュー}

AIに査読者やネイティブスピーカーの役割を担わせることで、論文の品質を事前に確認できる。

\begin{promptbox}
あなたは英語を母語とする[分野名]の研究者です。
以下の英語論文を読み、非ネイティブ話者が書いたと思われる
不自然な表現を指摘してください。

指摘の際は以下の形式でお願いします:
- 原文: [不自然な表現]
- 修正案: [自然な表現]
- 理由: [なぜ不自然か、どう改善されるか]

特に以下の点に注意してください:
1. 冠詞の使い方
2. 前置詞の選択
3. コロケーション（単語の自然な組み合わせ）
4. 文のリズムと読みやすさ

{[}英語論文を貼り付け]
\end{promptbox}

\begin{itembox}[l]{Tip: 推敲のタイミング}
推敲は時間をおいて行うのが効果的である。書いた直後ではなく、翌日に見直すと、問題点が見つかりやすくなる。AIを使った推敲も同様に、初回の英訳から時間をおいてから行うとよい。新鮮な目で見ることで、より多くの改善点を発見できる。
\end{itembox}

%======================================================================
\section{Abstract作成の技法}
\label{sec:4.5}
\index{Abstract@Abstract}

\subsection{Abstractの構造}

Abstractは論文全体の縮図であり、読者が論文を読むかどうかを決める最も重要な部分である。以下の5要素を含めるのが標準的である。

\begin{enumerate}
\item \textbf{背景（Background）:} 研究分野の文脈と解決すべき問題（1--2文）
\item \textbf{目的（Objective）:} 本研究で何を達成するか（1文）
\item \textbf{手法（Method）:} 提案するアプローチの概要（2--3文）
\item \textbf{結果（Results）:} 主要な実験結果の数値（1--2文）
\item \textbf{意義（Significance）:} 研究の意義と影響（1文）
\end{enumerate}

\subsection{150--200語にまとめるプロンプト設計}

\begin{promptbox}
以下の情報をもとに、学術論文のAbstractを英語で作成してください。

【研究情報】
- 研究分野: [分野]
- 解決する問題: [問題の記述]
- 提案手法: [手法の概要]
- 主要な結果: [数値を含む結果]
- 研究の意義: [意義の記述]

【条件】
- 語数: 150-200語
- 構造: 背景→目的→手法→結果→意義の順
- 時制: 背景は現在形、手法と結果は過去形、意義は現在形
- 具体的な数値を必ず含める
- 略語を初出で定義する
\end{promptbox}

\subsection{良いAbstractと悪いAbstractの比較}

\textbf{悪い例（問題点を含む）:}

\begin{quote}
In this paper, we propose a new method for text classification. Our method uses deep learning. We conducted experiments and obtained good results. Our method outperformed baseline methods. This work contributes to the field of natural language processing.
\end{quote}

\textbf{問題点:}
\begin{itemize}
\item 具体性がない（``new method''、``good results''）
\item 数値が一切含まれていない
\item 手法の特徴が不明
\item 読者にとって有益な情報がほとんどない
\end{itemize}

\textbf{良い例:}

\begin{quote}
Text classification in low-resource languages remains challenging due to the scarcity of labeled training data. We propose CrossLingual-BERT, a cross-lingual transfer learning framework that leverages pre-trained multilingual representations to classify texts in languages with limited annotated corpora. Our approach fine-tunes a multilingual BERT model on high-resource source languages and adapts it to target languages using only 100 labeled examples per class. Experiments on five typologically diverse languages demonstrate that CrossLingual-BERT achieves an average F1 score of 87.3\%, outperforming monolingual baselines by 12.5 percentage points and the previous state-of-the-art cross-lingual method by 4.2 percentage points. These results suggest that effective cross-lingual transfer is possible even with minimal target-language supervision, opening new possibilities for NLP applications in under-resourced languages.
\end{quote}

\textbf{優れている点:}
\begin{itemize}
\item 背景が具体的で問題が明確
\item 手法の核心的な特徴が述べられている
\item 具体的な数値が含まれている
\item 研究の意義が最後に述べられている
\end{itemize}

\begin{itembox}[l]{Tip: Abstractの執筆タイミング}
Abstractは論文の最後に書くのが効率的である。全セクションを書き終えてからAbstractを作成すれば、研究の全体像を正確に反映できる。AIにAbstractの草稿を作成させる際も、論文の完成後に全セクションの内容を入力として与えるとよい。
\end{itembox}

%======================================================================
\section{よくあるミスと対策}
\label{sec:4.6}

本節では、日本語話者が英語論文を書く際に陥りやすいミスとその対策を、授業で扱った具体例とともに体系的に解説する。まず各トピックの概念的な基礎を理解し、その後に早見表を参照するという流れで学習を進めてほしい。

\begin{screen}
\textbf{松野先生の教え：} 結局、日本語（母国語）の能力が重要である。必要最小限の日本語にしてから英語にする。これが技術英語の基本である。
\end{screen}

\subsection{すっきり最小主義：冗長な日本語を削ぎ落とす}
\index{すっきりさいしょうしゅぎ@すっきり最小主義}

英語に翻訳する前に、まず日本語を簡潔にすることが最も重要なステップである。冗長な日本語からは冗長な英語しか生まれない。

\subsubsection{重複表現を削る}

日本語は同じ意味の言葉を重ねる傾向がある。英訳前にこれを削ぎ落とす。

\begin{itembox}[l]{例：重複表現の除去}
\textbf{原文：} 「リチャージャブルバッテリーの安定した充電のためには、電気が流れる端子部が清浄であることが重要です」\\
\textbf{問題：} 「リチャージャブル（再充電可能な）」と「充電」が重複。「電気が流れる」と「端子」も重複。\\
\textbf{整理後：} 「安定した充電のためには端子を清潔に保つこと」\\
\textbf{英訳：} Clean the terminals for stable charging.
\end{itembox}

\subsubsection{空っぽ表現を捨てる}

日本語には意味のない「飾り」表現が多い。これを除去してから英訳する。

\begin{itembox}[l]{例：空っぽ表現の除去}
\textbf{原文：} 「オフィスでの集中力の妨げの原因となる不快な騒音を排除することで、生産性の向上が期待できるようになる」\\
\textbf{整理：} 「妨げの原因となる」→「妨げる」、「期待できるようになる」→不要\\
\textbf{整理後：} 「騒音を排除すれば生産性が向上する」\\
\textbf{英訳：} Eliminating noise improves productivity.
\end{itembox}

確信度を調整したい場合は、助動詞を使う。
\begin{itemize}
\item 弱い主張（5\%〜10\%程度）: Eliminating noise \textit{can} improve productivity.
\item 中程度の主張: Eliminating noise \textit{may} improve productivity.
\item 強い主張: Eliminating noise \textit{will} improve productivity.
\end{itemize}

\subsubsection{無駄な丁寧表現を除去する}

日本語の丁寧表現は英語に直訳すると不自然になる。

\begin{itembox}[l]{例：丁寧表現の除去}
\textbf{原文：} 「本日、無事に設計の最終確認を終了することができました」\\
\textbf{問題：} 「本日」「無事に」「終了することができました」は英語では不要\\
\textbf{整理後：} 「設計の最終確認を完了した」\\
\textbf{英訳：} We completed the final design review.
\end{itembox}

\begin{promptbox}
以下の日本語から、重複表現・空っぽ表現・無駄な丁寧表現を\\
取り除き、簡潔な日本語に整理してください。\\
その後、整理した日本語を英訳してください。\\
\\
{[}日本語原文を貼り付け{]}
\end{promptbox}

%---
\subsection{まず具体化主義：曖昧な表現を具体的にする}
\index{ぐたいかしゅぎ@具体化主義}

日本語の曖昧な表現は、英語にすると意味が通らなくなる。英訳前に具体的な表現に置き換える。

\begin{itembox}[l]{例：曖昧表現の具体化}
\textbf{原文：} 「オイル量のチェックはできるだけ頻繁に行うことが望ましい」\\
\textbf{問題：} 「できるだけ頻繁に」は英語にしにくい（as frequently as possible は不自然）\\
\textbf{具体化：} 「オイル量は1日1回点検すること」\\
\textbf{英訳：} Check the oil level once a day.
\end{itembox}

数値の曖昧さにも注意が必要である。

\begin{itembox}[l]{例：数値の具体化}
\textbf{原文：} 「生産は数千万だ」\\
\textbf{問題：} 「数千万」は英語でどう訳すか不明確\\
\textbf{具体化：} 「年間生産台数は3,600万台である」\\
\textbf{英訳：} Annual production is 36 million units.
\end{itembox}

略語も曖昧表現の一種である。初出時には必ずスペルアウトする。

\begin{itembox}[l]{例：略語のスペルアウト}
\textbf{推奨：} an antilock braking system (ABS)\\
\textbf{許容：} an ABS (antilock braking system)\\
\textbf{不可：} いきなり ABS と書く（読者が意味を知らない可能性がある）
\end{itembox}

%---
\subsection{とことん動詞主義：名詞句を動詞に戻す}
\index{どうししゅぎ@動詞主義}

日本語では「〜の実施」「〜の確認」「〜を行う」のように動詞を名詞化する傾向がある。これをそのまま英訳すると、名詞句が連鎖した読みにくい英文になる。動詞に戻してから英訳する。

\begin{itembox}[l]{例：名詞句 → 動詞}
\textbf{原文：} 「付録1はユニットの分解手順の説明である」\\
\textbf{名詞的英訳：} Appendix 1 is a description of the disassembly procedure of the unit.\\
\textbf{問題：} 前置詞 of が3回連続し、読みにくい\\
\textbf{動詞的英訳：} Appendix 1 describes how to disassemble the unit.
\end{itembox}

\begin{itembox}[l]{例：「行う」を削る}
\textbf{原文：} 「両企業はライセンス契約の修正を行った」\\
\textbf{名詞的英訳：} The two companies made an amendment to the license agreement.\\
\textbf{動詞的英訳：} The two companies amended the license agreement.
\end{itembox}

動詞を使うメリットは2つある。
\begin{enumerate}
\item \textbf{前置詞が減る：} 名詞句には前置詞（of, to, for等）が必要だが、動詞にすれば不要になることが多い。
\item \textbf{文が短くなる：} 「〜の実施を行った」（名詞+動詞）が「〜を実施した」（動詞のみ）になる。
\end{enumerate}

\begin{screen}
\textbf{4つの原則のまとめ：} 英訳の前に日本語を整える4つのステップ：（1）すっきり最小主義で冗長さを削り、（2）具体化主義で曖昧さをなくし、（3）動詞主義で名詞句を動詞に戻し、（4）能動態主義で受動態を能動態に変える。この4ステップを踏むだけで、AIによる英訳の品質は格段に上がる。
\end{screen}

%----------------------------------------------------------------------
\subsection{能動態と受動態の使い分け}
\label{subsec:voice}
\index{のうどうたい@能動態}
\index{じゅどうたい@受動態}

\subsubsection{「まっすぐ能動態主義」の原則}

日本語は受動態を多用する言語であるが、英語---特に技術英語や学術英語---では能動態が好まれる。「すっきり最小主義」（余計な情報を削ぎ落とす）と「まっすぐ能動態主義」（まず能動態で書くことを第一選択とする）を基本原則として覚えておくとよい。

日本語の技術文書では、「〜されている」「〜が行われた」といった受動態が自然であり、多用される。しかし、これをそのまま英語にすると冗長で読みにくい文章になる。英語では「誰が何をしたか」を明確にする能動態がストレートで好まれる。

\begin{itembox}[l]{例1: 日本語的な受動態から英語的な能動態へ}
\begin{itemize}
\item \textbf{日本語:} 「作業油の流量は自動で調整されている」
\item \textbf{受動態英訳:} ``The flow of the fluid is automatically regulated.''
\item \textbf{能動態英訳:} ``Pressure-sensitive valves automatically regulate the flow of the fluid.'' （主語が明確でストレート）
\end{itemize}
\end{itembox}

\begin{itembox}[l]{例2: 学術論文での能動態と受動態}
\begin{itemize}
\item \textbf{受動態:} ``The experiment was conducted by the authors.'' （冗長）
\item \textbf{能動態:} ``The authors conducted the experiment.'' （よりストレート）
\item \textbf{受動態:} ``It was observed that the accuracy was improved.'' （回りくどい）
\item \textbf{能動態:} ``We observed that the accuracy improved.'' （明確）
\end{itemize}
\end{itembox}

\subsubsection{受動態を使うべき4つの場面}

「まっすぐ能動態主義」とはいえ、受動態が適切な場面もある。以下の4つのケースでは、受動態を積極的に使うべきである。

\begin{enumerate}
\item \textbf{主語の統一（Subject Consistency）:} パラグラフ内で主語を統一するために受動態を使う。文ごとに主語がころころ変わると読みにくくなる。
\item \textbf{主語の短縮（Shorter Subject）:} 能動態にすると主語が非常に長くなる場合、受動態を使って主語を短くする。
\item \textbf{責任のあいまい化（Ambiguating Agency）:} 行為者を明示したくない場合、または行為者が重要でない場合に受動態を使う。
\item \textbf{話題の流れ（Topic Flow）:} 既知情報を文頭に、新情報を文末に配置するために受動態を使う。英語では「旧情報→新情報」の流れが自然である。
\end{enumerate}

\begin{itembox}[l]{受動態が適切な場面の例}
\begin{itemize}
\item 主語の統一: ``The samples were collected from three sites. They were then filtered and analyzed.'' （``The samples''を主語に統一）
\item 手法の説明: ``The samples were analyzed using mass spectrometry.'' （誰がやったかより何をしたかが重要）
\item 責任のあいまい化: ``Mistakes were made in the initial calibration.'' （誰の責任かを明示しない）
\end{itemize}
\end{itembox}

\begin{itembox}[l]{注意: 「うっかり受け身」は通じない}
日本語では「うっかり壊してしまった」のように、意図せずにそうなったことを受身で表現する（迷惑の受身）。しかし英語には同じ機能の受動態がないため、日本語の感覚でそのまま受動態を使うと意味が通じなくなる。英語では意図の有無にかかわらず、行為者を主語にして能動態で書くのが原則である。
\end{itembox}

\begin{promptbox}
以下の英文の受動態を能動態に書き換えてください。ただし、受動態の方が適切な場合は
その理由（主語の統一・主語の短縮・責任のあいまい化・話題の流れのいずれか）とともに
残してください。

{[}英文を貼り付け]
\end{promptbox}

以下に、受動態から能動態への代表的な改善パターンを早見表としてまとめる。

\begin{table}[htbp]
\centering
\caption{受動態から能動態への改善例}
\label{tab:passive}
\footnotesize
\begin{tabular}{ll}
\toprule
受動態（冗長） & 能動態（簡潔） \\
\midrule
The model was trained by us using... & We trained the model using... \\
It was observed that the accuracy & We observed that the accuracy \\
\quad was improved by... & \quad improved by... \\
The data was collected from... & We collected data from... \\
\bottomrule
\end{tabular}
\end{table}

\begin{itembox}[l]{Tip: ``We''の使用}
``We''を主語にすることに抵抗がある方もいるが、現代の学術英語では``We''の使用は一般的に受け入れられている。ただし、過度に``We''を繰り返さないよう、受動態と能動態を適切に混ぜることが理想的である。
\end{itembox}

%----------------------------------------------------------------------
\subsection{冠詞（a/the）の使い分け}
\label{subsec:articles}
\index{かんし@冠詞}

\subsubsection{なぜ冠詞が難しいのか}

冠詞（a, an, the, 無冠詞）は、日本語にない概念であるため、日本語話者にとって最も間違いやすい文法事項である。冠詞の選択は、「読者がその名詞をどの程度特定できるか」という情報の共有状態に基づいている。

\subsubsection{theの3つの用法}

定冠詞theには、大きく分けて3つの用法がある。

\begin{enumerate}
\item \textbf{唯一照応:} 世界に一つしかないもの、文脈上唯一に定まるもの。\\
例: ``the sun''（太陽は一つ）、``the Internet''（インターネットは一つ）、``the proposed method''（提案手法は一つ）
\item \textbf{前方照応:} 前の文で既に導入されたものを指す。a/anで初出した名詞を、2回目以降theで受ける。\\
例: ``Click the button to open \textbf{a} new window. Drag the images onto \textbf{the} new window.''（aで導入→theで再参照）
\item \textbf{後方照応:} 名詞の直後に修飾語句（関係節、前置詞句など）があり、それによって特定されるもの。\\
例: ``Locate \textbf{the} four screws that are behind the lid.''（``that are behind the lid''によって特定される）
\end{enumerate}

\begin{itembox}[l]{学術論文での冠詞の具体例}
\begin{itemize}
\item 初出: ``We propose \textbf{a} method for text classification.''
\item 再出: ``\textbf{The} method consists of three modules.''（前方照応）
\item 特定: ``\textbf{The} method proposed by Smith et al. (2023)''（後方照応）
\item 唯一: ``\textbf{The} proposed method outperforms all baselines.''（唯一照応）
\item よくある間違い: ``We propose method'' → ``We propose \textbf{a} method''（冠詞の欠落）
\end{itemize}
\end{itembox}

\subsubsection{可算名詞と不可算名詞}

同じ単語でも、可算名詞と不可算名詞で意味が変わることがある。冠詞の有無が意味を決定するため、注意が必要である。

\begin{itemize}
\item \textbf{iron}（不可算）= 鉄（物質） / \textbf{an iron}（可算）= アイロン（道具）
\item \textbf{chicken}（不可算）= 鶏肉（食材） / \textbf{a chicken}（可算）= ニワトリ（動物）
\item \textbf{paper}（不可算）= 紙（素材） / \textbf{a paper}（可算）= 論文
\item \textbf{glass}（不可算）= ガラス（素材） / \textbf{a glass}（可算）= コップ
\end{itemize}

\begin{promptbox}
以下の英文の冠詞（a/an/the）の使い方を確認し、誤りがあれば修正してください。
修正した場合は、唯一照応・前方照応・後方照応のどれに該当するか説明してください。

{[}英文を貼り付け]
\end{promptbox}

以下に、冠詞の基本ルールを早見表としてまとめる。

\textbf{基本ルール:}
\begin{itemize}
\item \textbf{the:} 特定のもの、前出のもの、唯一のもの（唯一照応・前方照応・後方照応）
\item \textbf{a/an:} 不特定のもの、初出のもの
\item \textbf{無冠詞:} 複数形の一般論、不可算名詞の一般論
\end{itemize}

\begin{promptbox}
以下の英語論文の冠詞の使用を確認してください。
特に以下のパターンに注意してチェックしてください:

1. 初出の名詞に "the" を使っていないか
2. 前出の名詞に "a/an" を使っていないか
3. 数式や手法名に不要な冠詞がついていないか
4. "the proposed method" のように特定される名詞に
   "the" が付いているか

{[}英語論文を貼り付け]
\end{promptbox}

%----------------------------------------------------------------------
\subsection{時制の選択}
\label{subsec:tense}
\index{じせい@時制}

\subsubsection{論文のセクションと時制の関係}

英語論文では、セクションによって適切な時制が決まっている。これは単なる慣習ではなく、各セクションで述べる内容の性質に基づいている。

\begin{itemize}
\item \textbf{Methods（手法）:} \textbf{過去形}を使う。実験は過去に行ったことだからである。\\
例: ``We trained the model using 10,000 samples.''
\item \textbf{Results（結果）:} \textbf{過去形}を使う。結果も過去に得たものだからである。\\
例: ``The proposed method achieved an accuracy of 95.3\%.''
\item \textbf{Introduction / Discussion:} \textbf{現在形}を使う。一般的な事実や先行研究の知見は現在も成り立つものだからである。\\
例: ``Deep learning has become a dominant approach in NLP.''
\end{itemize}

\subsubsection{単純現在形の本質: 再現性100\%}

単純現在形は、「いつでも・誰がやっても成り立つ事実」を述べるときに使う。つまり、再現性が100\%であることを主張する時制である。

\begin{itemize}
\item ``Water boils at 100°C.''（水は100°Cで沸騰する = いつでも成り立つ事実）
\item ``This algorithm converges in $O(n \log n)$ time.''（このアルゴリズムは$O(n \log n)$で収束する = 数学的事実）
\end{itemize}

逆に言えば、単純現在形を使うということは「これは普遍的な事実である」と主張していることになる。自分の実験結果に単純現在形を使うと、「この結果はいつでも再現される」という強い主張になるので注意が必要である。

\subsubsection{現在完了形の3つの用法}

現在完了形は学術論文で頻繁に使われるが、以下の3つの用法を区別して使う必要がある。

\begin{enumerate}
\item \textbf{完了（already）:} 動作が完了したことを表す。\\
例: ``We have completed the analysis.''（分析を完了した）
\item \textbf{経験（times）:} 過去の経験を表す。\\
例: ``Several studies have shown that this approach is effective.''（複数の研究が示してきた）
\item \textbf{継続（for / since）:} 過去から現在まで続いていることを表す。\\
例: ``This method has been used for over a decade.''（10年以上使われ続けている）
\end{enumerate}

\subsubsection{助動詞の使い分け}

助動詞は、話者の「確信度」を表す重要な手段である。学術論文では、主張の強さを適切にコントロールするために助動詞を使い分ける。

\begin{itemize}
\item \textbf{will} = 確信・コミットメント（確信度が高い）\\
例: ``This approach will enable faster training.''
\item \textbf{can} = 可能性・能力\\
例: ``This method can handle large-scale datasets.''
\item \textbf{may} = 不確実性・推測（確信度が低い）\\
例: ``This discrepancy may be due to overfitting.''
\end{itemize}

\begin{itembox}[l]{注意: ``cannot'' と ``can not'' の違い}
``cannot''（1語）は「〜できない」という不可能の意味である。一方、``can not''（2語）は``can not only A but also B''のように、``not''が後続の語句を否定する場合に使う。通常の「〜できない」を表すときは必ず``cannot''と1語で書く。
\end{itembox}

以下に、セクション別の適切な時制を早見表としてまとめる。

\begin{table}[htbp]
\centering
\caption{セクション別の適切な時制}
\label{tab:tense}
\footnotesize
\begin{tabular}{lll}
\toprule
セクション & 基本の時制 & 例 \\
\midrule
Introduction & 現在形（一般的事実）/ & ``Deep learning has achieved...'' \\
 & 過去形（先行研究の報告） & ``Smith et al. (2023) proposed...'' \\
Methods & 過去形 & ``We trained the model using...'' \\
Results & 過去形 & ``The proposed method achieved...'' \\
Discussion & 現在形（解釈）/ & ``These results suggest...'' \\
 & 過去形（結果への言及） & ``Our model showed...'' \\
\bottomrule
\end{tabular}
\end{table}

\begin{promptbox}
以下の英語論文のMethodsセクションについて、
時制が適切に使われているか確認してください。
Methodsでは基本的に過去形を使用します。
不適切な時制が使われている箇所を指摘し、修正してください。

{[}Methodsセクションを貼り付け]
\end{promptbox}

%----------------------------------------------------------------------
\subsection{自動詞と他動詞の注意}
\label{subsec:transitive}
\index{じどうし@自動詞}
\index{たどうし@他動詞}

英語では、同じ動詞でも自動詞として使うか他動詞として使うかで意味が大きく変わることがある。日本語では助詞（「を」「に」など）で関係を示すため、動詞の自他をあまり意識しないが、英語では目的語の有無が意味を決定する。

\begin{itembox}[l]{自動詞と他動詞で意味が変わる例}
\begin{itemize}
\item \textbf{submit}
  \begin{itemize}
  \item 他動詞: ``We submitted the paper.''（論文を提出した）
  \item 自動詞: ``The army submitted.''（軍が降伏した・屈服した）
  \end{itemize}
\item \textbf{search}
  \begin{itemize}
  \item 他動詞: ``I searched him.''（彼をボディチェックした = 身体検査した）
  \item 自動詞+前置詞: ``I searched for him.''（彼を探した）
  \end{itemize}
\item \textbf{run}
  \begin{itemize}
  \item 自動詞: ``The program runs.''（プログラムが動作する）
  \item 他動詞: ``We run the program.''（プログラムを実行する）
  \end{itemize}
\end{itemize}
\end{itembox}

特に``search''は要注意である。``search the database''は「データベースの中を探す」（正しい用法）だが、``search him''は「彼の身体を捜索する」という意味になり、「彼を探す」とは全く異なる。「〜を探す」と言いたい場合は``search for''を使う。

学術論文で動詞を使う際は、辞書で自動詞・他動詞の区別を必ず確認すること。特に、日本語で「〜を」と言える動詞が英語でも他動詞とは限らない点に注意が必要である。

%----------------------------------------------------------------------
\subsection{直訳による不自然な表現}
\index{ちょくやく@直訳}

日本語を直訳すると、文法的には正しくても不自然な英語になることがある。

\begin{table}[htbp]
\centering
\caption{直訳による不自然な表現と自然な英語}
\label{tab:unnatural}
\footnotesize
\begin{tabular}{ll}
\toprule
日本語直訳（不自然） & 自然な英語 \\
\midrule
``We could confirm that...'' & ``We confirmed that...'' \\
``It is considered that...'' & ``We consider...'' / ``This suggests...'' \\
``From this result...'' & ``Based on these results...'' \\
``We performed the experiment...'' & ``We conducted experiments...'' \\
``The accuracy became high...'' & ``The accuracy improved...'' \\
``As a merit of our method...'' & ``An advantage of our method is...'' \\
\bottomrule
\end{tabular}
\end{table}

\begin{promptbox}
以下の英語論文について、日本語の直訳に起因する不自然な表現を
特定してください。各箇所について、より自然な英語表現を提案し、
なぜその表現の方が自然かを説明してください。

{[}英語論文を貼り付け]
\end{promptbox}

%----------------------------------------------------------------------
\subsection{専門用語の確認方法}
\index{せんもんようご@専門用語}

専門用語の英語表現が分からない場合や、正しいかどうか確認したい場合は、以下の方法が有効である。

\begin{enumerate}
\item \textbf{Google Scholarでの用例検索:} ``用語''を含む論文を検索し、トップ論文でどのように使われているか確認する
\item \textbf{AIに確認を依頼:}
\end{enumerate}

\begin{promptbox}
「[日本語の専門用語]」の英語表現として、以下のどちらが
{[}分野名]の学術論文で一般的ですか？

A: [候補1]
B: [候補2]

それぞれの使用頻度や文脈の違いがあれば教えてください。
また、他に一般的な表現があれば提示してください。
\end{promptbox}

%======================================================================
\section{学術英語表現集}
\label{sec:4.7}
\index{がくじゅつえいごひょうげん@学術英語表現}

\subsection{Introductionの定型表現}

\begin{table}[htbp]
\centering
\caption{Introductionで使える定型表現}
\label{tab:intro-expressions}
\footnotesize
\begin{tabular}{lp{0.55\textwidth}}
\toprule
目的 & 表現例 \\
\midrule
分野の重要性を述べる & ``X has attracted significant attention in recent years.'' \\
 & ``X plays a crucial role in...'' \\
 & ``X is of fundamental importance to...'' \\
\midrule
先行研究を紹介する & ``Several studies have investigated...'' \\
 & ``Previous work has shown that...'' \\
 & ``Smith et al. (2023) demonstrated that...'' \\
\midrule
問題点を指摘する & ``However, existing methods suffer from...'' \\
 & ``Despite these advances, ... remains challenging.'' \\
 & ``A major limitation of prior work is...'' \\
\midrule
本研究の目的を述べる & ``In this paper, we propose...'' \\
 & ``This work addresses the problem of...'' \\
 & ``We present a novel approach to...'' \\
\midrule
貢献を列挙する & ``Our main contributions are as follows:'' \\
 & ``The key contributions of this work include:'' \\
\bottomrule
\end{tabular}
\end{table}

\subsection{Methodsの定型表現}

\begin{table}[htbp]
\centering
\caption{Methodsで使える定型表現}
\label{tab:methods-expressions}
\footnotesize
\begin{tabular}{lp{0.55\textwidth}}
\toprule
目的 & 表現例 \\
\midrule
手法の全体像を述べる & ``Our approach consists of three main components:'' \\
 & ``The overall architecture is illustrated in Figure X.'' \\
 & ``We formulate the problem as...'' \\
\midrule
手順を説明する & ``First, we... Then, we... Finally, we...'' \\
 & ``The input is processed through...'' \\
 & ``We apply X to obtain Y.'' \\
\midrule
数式を導入する & ``The objective function is defined as:'' \\
 & ``We optimize the following loss function:'' \\
 & ``where X denotes the...'' \\
\midrule
実装の詳細を述べる & ``We implemented our method using...'' \\
 & ``All experiments were conducted on...'' \\
 & ``The hyperparameters were tuned via...'' \\
\bottomrule
\end{tabular}
\end{table}

\subsection{Resultsの定型表現}

\begin{table}[htbp]
\centering
\caption{Resultsで使える定型表現}
\label{tab:results-expressions}
\footnotesize
\begin{tabular}{lp{0.55\textwidth}}
\toprule
目的 & 表現例 \\
\midrule
結果を提示する & ``Table X summarizes the results of...'' \\
 & ``As shown in Figure X, our method...'' \\
 & ``The results are presented in Table X.'' \\
\midrule
比較する & ``Our method outperforms the baseline by X\%.'' \\
 & ``Compared to the state-of-the-art, our approach achieves...'' \\
 & ``The improvement is statistically significant ($p < 0.05$).'' \\
\midrule
傾向を述べる & ``We observe a consistent improvement across...'' \\
 & ``The performance increases as X grows.'' \\
 & ``Notably, our method shows robust performance on...'' \\
\bottomrule
\end{tabular}
\end{table}

\subsection{Discussion / Conclusionの定型表現}

\begin{table}[htbp]
\centering
\caption{Discussion / Conclusionで使える定型表現}
\label{tab:discussion-expressions}
\footnotesize
\begin{tabular}{lp{0.55\textwidth}}
\toprule
目的 & 表現例 \\
\midrule
結果を解釈する & ``These results suggest that...'' \\
 & ``One possible explanation is that...'' \\
 & ``This finding is consistent with...'' \\
\midrule
限界を述べる & ``A limitation of our study is...'' \\
 & ``Our approach may not generalize to...'' \\
 & ``One potential concern is...'' \\
\midrule
将来の方向を述べる & ``Future work could explore...'' \\
 & ``A promising direction is to...'' \\
 & ``We plan to extend this work by...'' \\
\midrule
結論を述べる & ``In summary, we have presented...'' \\
 & ``This work demonstrates the effectiveness of...'' \\
 & ``Our findings highlight the potential of...'' \\
\bottomrule
\end{tabular}
\end{table}

\begin{itembox}[l]{Tip: 定型表現の活用法}
これらの表現を暗記する必要はない。執筆時に参照表として使い、自分の分野の論文でよく見かける表現を徐々に身につけていくとよい。定型表現は「型」として活用し、その中に自分の研究の独自の内容を盛り込むことが重要である。
\end{itembox}

%======================================================================
\section*{演習問題}

\begin{enumerate}

\item \textbf{演習1: IMRaD構造の分析}

自分の研究分野のトップ会議またはジャーナルから論文を1本選び、以下を行いなさい。
\begin{enumerate}
\item 各セクション（Introduction, Methods, Results, Discussion）の段落数と概算語数を調べる
\item Introductionの「逆三角形構造」（背景→問題→本研究の貢献）を具体的に特定する
\item 各段落のトピックセンテンスを抜き出し、それだけを読んで論文の流れが理解できるか確認する
\item この分析から得られた、英語論文の構造に関する気づきをまとめる
\end{enumerate}

\item \textbf{演習2: AIを使った英訳の実践}

自分の研究に関する日本語の文章（300--500字）を用意し、以下を行いなさい。
\begin{enumerate}
\item 本章のStep 1のチェックリストを使って日本語原稿を確認・修正する
\item Step 2のプロンプトテンプレートを使って、AIに英訳を依頼する
\item Step 3のプロンプトを使って、英訳を学術的に洗練する
\item Step 2の英訳とStep 3の洗練後の英訳を比較し、どのような改善がなされたかを分析する
\item 自分の分野の論文と表現を比較し、さらなる改善点がないか検討する
\end{enumerate}

\item \textbf{演習3: 逆翻訳チェックの実践}

演習2で作成した英語論文の一部について、以下を行いなさい。
\begin{enumerate}
\item 別のAIツール（演習2でClaudeを使った場合はChatGPT等）で日本語に逆翻訳する
\item 逆翻訳された日本語を元の日本語原稿と比較し、意味のズレを特定する
\item ズレが見つかった箇所について、英訳のどこに問題があるかを分析する
\item 英訳を修正し、再度逆翻訳して意味のズレが解消されたか確認する
\end{enumerate}

\item \textbf{演習4: Abstract作成の実践}

自分の研究（または研究計画）について、以下を行いなさい。
\begin{enumerate}
\item 本章のAbstract構造（背景→目的→手法→結果→意義）に従い、各要素を日本語で箇条書きにする
\item 本章のプロンプトテンプレートを使って、AIに英語のAbstractを作成させる
\item 作成されたAbstractが150--200語に収まっているか確認する
\item 本章の「良いAbstractの条件」（具体性、数値、構造）を満たしているか評価する
\item 必要に応じて修正し、最終版を完成させる
\end{enumerate}

\item \textbf{演習5: 総合推敲の実践}

演習2--4で作成した英語テキストをまとめ、以下の段階的推敲を行いなさい。
\begin{enumerate}
\item 第1段階として、文法チェックのプロンプトを使って文法エラーを修正する
\item 第2段階として、表現チェックのプロンプトを使って学術的表現に洗練する
\item 第3段階として、論理チェックのプロンプトを使って論理構造を確認する
\item 最後に、ネイティブチェック風レビューのプロンプトを使って自然さを確認する
\item 各段階で見つかった問題の種類と数を記録し、自分の英語ライティングの傾向を分析する
\end{enumerate}

\end{enumerate}

\vspace{10mm}
\begin{itembox}[l]{本章のまとめ}
英語論文の執筆は、日本語話者にとって大きな挑戦であるが、AIを活用することでそのハードルを大幅に下げることができる。重要なのは、AIを「翻訳機」としてではなく「執筆パートナー」として活用することである。AIの出力をそのまま使うのではなく、段階的な推敲プロセスを通じて品質を高め、最終的には自分の言葉として責任を持てる英語論文に仕上げることが大切である。繰り返し実践することで、AIに頼らなくても書ける力が自然と身についていく。
\end{itembox}
